The hit rate is low. RBX1 was picked to be hard, and 2.8\% reflects
that. The two prior Adaptyv community campaigns (EGFR 2024 at
\textasciitilde21\%, Nipah-G 2025--26 at \textasciitilde9\% on the
public collection) used larger, more druggable targets with deeper
prior art. RBX1 has neither. The disordered N-terminal tail offers no
stable interface. The RING core depends on metal coordination that the
default-buffer screen does not enforce.

The methods that produced binders do not form a clean ranking.
RFdiffusion sent the most submissions and placed no binder. ORBIT
(Mandrake Bio, unpublished) placed the Strong hit with a single 100-aa
miniprotein. BindCraft 2 {[}@pacesa2025bindcraft{]} produced three
Medium binders from seven submissions, a 43\% hit rate on a slice too
thin for a significance claim. AF2 hallucination with ADFlip produced
the next two tightest. Four distinct design philosophies (full
diffusion, AF2 backprop with a sequence editor, composite-objective
JAX, inverse-folding-style) each delivered a top-five binder.

Mini binders dominate the binder set. The submission pool was already
25\% miniprotein, 13\% antibody-family, 10\% peptide, and 51\%
\emph{Other}, because most publicly available design tools target
miniproteins. The binder-set bias likely reflects the input bias, not
an intrinsic advantage at this target.

The 322 sequences, 322 ESMFold predictions, and 943 SPR replicate
traces are a public benchmark for model calibration on a hard target.
The next follow-up is a Zn²⁺-loaded rerun, which probes whether
designs are targeting the apo or holo conformation.

RBX1 is not yet a designable target. One binder sits under 100 nM, two
between 100 and 200 nM, and six between 200 and 500 nM. Design
pipelines have not generalised to small, metal-coordinated, partially
disordered targets. The 322-design dataset shows how thin the binder
pool is at this difficulty and where the next round of follow-ups
should land.
